\documentclass[11pt,a4paper]{article}

\usepackage[T1]{fontenc}
\usepackage[utf8]{inputenc}
\usepackage{lmodern}
\usepackage{microtype}
\usepackage{geometry}
\geometry{top=2.6cm, bottom=2.6cm, left=3cm, right=3cm}

\usepackage{amsmath}
\usepackage{amssymb}
\usepackage{graphicx}
\usepackage{booktabs}
\usepackage{xcolor}
\usepackage{titlesec}
\usepackage{parskip}
\usepackage{enumitem}
\usepackage{fancyhdr}
\usepackage[numbers,sort&compress]{natbib}
\usepackage[hidelinks]{hyperref}

%% Colours
\definecolor{camblue}{HTML}{003A6C}
\definecolor{rulecol}{HTML}{003A6C}
\definecolor{subtlegray}{HTML}{555555}

%% Section headings: bold, small-caps feel, with a thin rule
\titleformat{\section}
  {\normalfont\normalsize\bfseries\color{camblue}}
  {}{0em}
  {}
  [\vspace{1pt}{\color{rulecol}\rule{\linewidth}{0.5pt}}\vspace{-2pt}]
\titlespacing{\section}{0pt}{14pt plus 2pt minus 2pt}{6pt plus 1pt}

%% Footer only
\pagestyle{fancy}
\fancyhf{}
\renewcommand{\headrulewidth}{0pt}
\fancyfoot[R]{\small\color{subtlegray}\thepage}
\fancyfoot[L]{\small\color{subtlegray}Executive Summary}

%% Tighten list spacing
\setlist[itemize]{nosep, leftmargin=1.4em, topsep=4pt}

\begin{document}

%% ---------------------------------------------------------------
%%  HEADER
%% ---------------------------------------------------------------
\thispagestyle{fancy}

\noindent
{\color{rulecol}\rule{\linewidth}{1.8pt}}\\[4pt]
{\LARGE\bfseries The Geometry of Concept Representation\\[2pt]
in Open-Source Language Models}\\[6pt]
{\color{rulecol}\rule{\linewidth}{0.5pt}}

\vspace{8pt}
\noindent
{\normalsize
\textbf{Elisabeta-Iulia (Julia) Dima}\enspace$\cdot$\enspace
MPhil Data Intensive Science\enspace$\cdot$\enspace
University of Cambridge, June 2026
}\\[2pt]
\noindent
{\small\color{subtlegray}
Supervised by Dr Miles Cranmer and Dr Alessandro Favero\enspace$\cdot$\enspace
Department of Applied Mathematics and Theoretical Physics
}

\vspace{14pt}

%% ---------------------------------------------------------------
\section{Motivation and Scientific Justification}
%% ---------------------------------------------------------------

Large language models (LLMs) routinely produce correct answers to arithmetic and reasoning problems, yet output accuracy cannot establish \emph{how} those answers are produced. A model can score highly on addition by exploiting digit-level heuristics or memorised templates rather than by computing the intended function. This distinction is not merely academic: if correct outputs rely on surface shortcuts, those shortcuts may fail under distributional shift or adversarial prompting, and no downstream audit based on accuracy alone would detect the failure.

Mechanistic interpretability addresses this by reverse-engineering the internal computations of neural networks. Anthropic's recent paper \textit{On the Biology of a Large Language Model}~\cite{lindsey2025biology} demonstrated that sparse, interpretable circuits can be identified within large base models. Two significant gaps remain open, however. First, it is unknown whether analogous circuits arise in instruction-tuned models — which are the models actually deployed. Second, existing techniques require manual, task-by-task circuit tracing; there is no general method for determining where an arbitrary abstract concept first emerges in the residual stream, or what its representational geometry looks like.

This thesis addresses both gaps through two complementary contributions: (1) a reproduction of Anthropic's circuit-tracing methodology on Qwen3-4B, an instruction-tuned 4-billion-parameter model, and (2) a new \emph{contrastive residual-stream delta} method for localising abstract concepts within the residual stream and characterising their geometry.

%% ---------------------------------------------------------------
\section{Methodology}
%% ---------------------------------------------------------------

\textbf{Concept localisation.}\enspace
For each concept under study, matched prompt pairs are constructed: one prompt instantiates the concept; the other does not; syntax, length, and answer format are held constant. The mean difference in residual-stream activations across layers and token positions — the \emph{delta trajectory} — reveals where the concept becomes geometrically salient. Trajectories are normalised against a permutation null baseline, removing confounds from layer-wise activation scale growth.

\emph{Anchor} positions (layer and token where the delta is large and non-monotone) are ranked by a non-monotonicity score. Causal sufficiency of each anchor is then tested by activation patching: replacing residual-stream vectors in negative-concept examples with those from positive examples and measuring the resulting change in the model's first output token.

\textbf{Feature attribution.}\enspace
At each anchor, the delta vector is projected through the layer's \emph{transcoder} — a sparse autoencoder trained to reconstruct each MLP output — to retrieve individual features that activate contrastively and write along the concept direction. Their joint geometry defines a candidate concept manifold; attribution-graph connections among them form a \emph{feature constellation} describing how concept-relevant information is routed and transformed.

\textbf{Tasks.}\enspace
The method is applied to three arithmetic tasks with Qwen3-4B, ordered by representational complexity:
\begin{itemize}
  \item \textbf{Carry detection} — does the ones-digit sum in a single-step addition carry?
  \item \textbf{GCD divisibility} — is $\gcd(a,7)=7$?
  \item \textbf{Residue class} — is $a \equiv 1\pmod{7}$?
\end{itemize}
A soft-prompt (learned prefix) experiment serves as a diagnostic control, testing whether input-space perturbations can redirect concept-relevant computation.

%% ---------------------------------------------------------------
\section{Key Findings}
%% ---------------------------------------------------------------

\textbf{Carry} exhibits the clearest and strongest geometric structure. A stable concept direction emerges at layers 4--5 at the final-operand token. The inter-layer cosine-similarity matrix shows a single high-similarity block, indicating the direction is \emph{maintained} rather than rotating across layers. Activation patching from positive to negative examples restores near-perfect accuracy, confirming causal sufficiency. Transcoder analysis at this anchor recovers features with structured activation patterns over the operand grid — iso-sum detectors, parity detectors, operand-comparison features — consistent with a Fourier-like arithmetic substrate. Ablating all Fourier-structured features jointly leaves accuracy unchanged, indicating that this feature set is geometrically present but not part of the causal bottleneck under the tested intervention.

\textbf{GCD divisibility} shows a qualitatively different structure: the inter-layer cosine matrix exhibits a sharp block transition at a specific layer, indicating that the concept direction \emph{rotates} rather than persists. Concept-relevant features are localisable but more distributed across transcoder feature space than for carry.

\textbf{Residue class} produces no stable direction above the null permutation band at any tested anchor or layer. This is consistent with the model's low first-token accuracy (12\%) and suggests the model does not linearly encode modular residue membership in a form accessible to contrastive analysis — a negative result that itself characterises the limits of the model's arithmetic representations.

\textbf{Reproducibility.}\enspace
Anthropic's circuit-tracing methodology extends to Qwen3-4B. Interpretable circuits for addition are recovered, including carry-routing components analogous to those described for base models, with differences in layer distribution attributable to instruction-tuning.

\textbf{Soft prompting.}\enspace
Learned prefixes can repair response-format failures: first-token accuracy rises from 0\% to 28.6\% on balanced parentheses, and from 0\% to 53.8\% on causal-direction tasks. Residue class accuracy increases by only 1 percentage point, and the mean correct-token probability \emph{decreases}. Learned prefix embeddings do not align with identified concept directions and are not interpretable in transcoder feature space — consistent with encoding a continuous approximation of the information the model requires rather than structured semantic content.

\begin{table}[h]
\centering
\small
\setlength{\tabcolsep}{8pt}
\begin{tabular}{@{}lrrrr@{}}
\toprule
Task & Base acc.\ (\%) & Prefix acc.\ (\%) & $\Delta$ (\%) & $\Delta\bar{p}(\text{correct})$ \\
\midrule
Residue class        & 15.0 & 16.0 & $+1.0$  & $-0.011$ \\
Balanced parentheses &  0.0 & 28.6 & $+28.6$ & $+0.078$ \\
Causal direction     &  0.0 & 53.8 & $+53.8$ & $+0.236$ \\
\bottomrule
\end{tabular}
\caption{\small Soft-prompt results for Qwen3-4B ($k=10$ prefix tokens, base weights frozen).}
\label{tab:sp_summary}
\end{table}

%% ---------------------------------------------------------------
\section{Recommendations and Next Steps}
%% ---------------------------------------------------------------

The contrastive delta method is designed to generalise: any domain where a binary concept can be expressed through matched prompt pairs is accessible. Immediate extensions include chemical reactivity, logical entailment, and spectral classification, where contrastive prompt pairs can be constructed from existing scientific datasets.

The Gram-matrix structure of the transcoder decoder projection — showing that cosine alignment scores conflate direct feature contributions with cross-feature interference — motivates a refined feature ranking criterion that accounts for co-activation geometry. This would improve the precision of feature constellation construction.

Combining delta localisation with full attribution-graph analysis would link the geometric picture to causal pathways, enabling a mechanistic account of \emph{why} certain concepts resist contrastive decomposition. Finally, applying the method across model scales would test whether the observed heterogeneity in concept geometry scales with model capacity.

%% ---------------------------------------------------------------
\section{Research Impact and Conclusion}
%% ---------------------------------------------------------------

This work delivers a model-agnostic concept localisation method that complements existing circuit-tracing approaches, and applies it to produce the first systematic geometric characterisation of arithmetic representations in an instruction-tuned LLM. The central finding — that concept geometry is non-uniform — has direct practical implications. Carry detection crystallises into a low-rank, causally active direction accessible to intervention. Residue class membership does not. This heterogeneity means that the success of activation-based steering, knowledge editing, or representation-based probing depends fundamentally on the representational form of the target concept, and that localisation should precede any editing attempt.

As LLMs are deployed in high-stakes scientific and decision-making contexts, tools that distinguish genuine internal computation from surface-level pattern matching become essential for building reliable and auditable AI systems. The pipeline developed here — contrastive delta trajectories, anchor-conditioned transcoder projections, and attribution-graph feature constellations — provides a reusable, open-source foundation for this kind of analysis.

%% ---------------------------------------------------------------
\vspace{8pt}
{\color{rulecol}\rule{\linewidth}{0.5pt}}

\begingroup
\small
\bibliographystyle{unsrt}
\bibliography{MPhilThesisReport/thesis}
\endgroup

\end{document}
